% ===== Graphs (13 problems) =====
\section{Graphs \texorpdfstring{{\normalsize\color{muted}$\cdot$ 13 problems}}{- 13 problems}}
\markboth{Graphs}{}

\begin{problem}{200}{Number of Islands}{Medium}
\summary{Count groups of connected '1's (4-directional) in a grid.}
\keyidea{Scan; at each unvisited land cell, count +1 and flood-fill (DFS/BFS) the whole island.}
\complexity{O(m * n) time and space}
\begin{code}
class Solution:
    def numIslands(self, grid: List[List[str]]) -> int:
        R, C = len(grid), len(grid[0])
        def sink(r, c):
            if 0 <= r < R and 0 <= c < C and grid[r][c] == "1":
                grid[r][c] = "0"
                sink(r + 1, c); sink(r - 1, c); sink(r, c + 1); sink(r, c - 1)
        count = 0
        for r in range(R):
            for c in range(C):
                if grid[r][c] == "1":
                    count += 1
                    sink(r, c)
        return count
\end{code}
\end{problem}

\begin{problem}{133}{Clone Graph}{Medium}
\summary{Deep-copy a connected undirected graph given one node.}
\keyidea{DFS with a map old -\textgreater{} clone; create the clone before visiting neighbours (handles cycles).}
\complexity{O(V + E) time, O(V) space}
\begin{code}
class Solution:
    def cloneGraph(self, node: Optional["Node"]) -> Optional["Node"]:
        clones = {}
        def dfs(n):
            if n in clones:
                return clones[n]
            copy = Node(n.val)
            clones[n] = copy
            for nb in n.neighbors:
                copy.neighbors.append(dfs(nb))
            return copy
        return dfs(node) if node else None
\end{code}
\end{problem}

\begin{problem}{695}{Max Area of Island}{Medium}
\summary{Size of the largest island of 1s.}
\keyidea{Flood-fill returning the area; take the max.}
\complexity{O(m * n) time and space}
\begin{code}
class Solution:
    def maxAreaOfIsland(self, grid: List[List[int]]) -> int:
        R, C = len(grid), len(grid[0])
        def area(r, c):
            if not (0 <= r < R and 0 <= c < C) or grid[r][c] != 1:
                return 0
            grid[r][c] = 0
            return 1 + area(r + 1, c) + area(r - 1, c) + area(r, c + 1) + area(r, c - 1)
        return max(area(r, c) for r in range(R) for c in range(C))
\end{code}
\end{problem}

\begin{problem}{417}{Pacific Atlantic Water Flow}{Medium}
\summary{Cells whose water can flow (to equal or lower neighbours) to both the Pacific (top/left) and Atlantic (bottom/right).}
\keyidea{Reverse it: DFS uphill from each ocean's border; answer = intersection of both reachable sets.}
\complexity{O(m * n) time and space}
\begin{code}
class Solution:
    def pacificAtlantic(self, heights: List[List[int]]) -> List[List[int]]:
        R, C = len(heights), len(heights[0])
        pac, atl = set(), set()
        def dfs(r, c, seen, prev):
            if ((r, c) in seen or not (0 <= r < R and 0 <= c < C)
                    or heights[r][c] < prev):
                return
            seen.add((r, c))
            for dr, dc in ((1, 0), (-1, 0), (0, 1), (0, -1)):
                dfs(r + dr, c + dc, seen, heights[r][c])
        for c in range(C):
            dfs(0, c, pac, heights[0][c])
            dfs(R - 1, c, atl, heights[R - 1][c])
        for r in range(R):
            dfs(r, 0, pac, heights[r][0])
            dfs(r, C - 1, atl, heights[r][C - 1])
        return [list(p) for p in pac & atl]
\end{code}
\end{problem}

\begin{problem}{130}{Surrounded Regions}{Medium}
\summary{Flip every 'O' region not connected to the border into 'X'.}
\keyidea{Mark border-connected O's as safe (T) by DFS; flip the rest to X; restore T to O.}
\complexity{O(m * n) time and space}
\begin{code}
class Solution:
    def solve(self, board: List[List[str]]) -> None:
        R, C = len(board), len(board[0])
        def mark(r, c):
            if 0 <= r < R and 0 <= c < C and board[r][c] == "O":
                board[r][c] = "T"
                mark(r + 1, c); mark(r - 1, c); mark(r, c + 1); mark(r, c - 1)
        for r in range(R):
            mark(r, 0); mark(r, C - 1)
        for c in range(C):
            mark(0, c); mark(R - 1, c)
        for r in range(R):
            for c in range(C):
                if board[r][c] == "O":
                    board[r][c] = "X"
                elif board[r][c] == "T":
                    board[r][c] = "O"
\end{code}
\end{problem}

\begin{problem}{994}{Rotting Oranges}{Medium}
\summary{Each minute rot spreads to adjacent fresh oranges. Minutes until none are fresh (-1 if impossible).}
\keyidea{Multi-source BFS from all rotten oranges at once; count levels.}
\complexity{O(m * n) time and space}
\begin{code}
class Solution:
    def orangesRotting(self, grid: List[List[int]]) -> int:
        R, C = len(grid), len(grid[0])
        q, fresh = deque(), 0
        for r in range(R):
            for c in range(C):
                if grid[r][c] == 2: q.append((r, c))
                elif grid[r][c] == 1: fresh += 1
        minutes = 0
        while q and fresh:
            for _ in range(len(q)):
                r, c = q.popleft()
                for nr, nc in ((r + 1, c), (r - 1, c), (r, c + 1), (r, c - 1)):
                    if 0 <= nr < R and 0 <= nc < C and grid[nr][nc] == 1:
                        grid[nr][nc] = 2
                        fresh -= 1
                        q.append((nr, nc))
            minutes += 1
        return minutes if fresh == 0 else -1
\end{code}
\end{problem}

\begin{problem}{286}{Walls and Gates}{Medium}
\summary{Fill each empty room (INF) with its distance to the nearest gate (0); walls are -1.}
\keyidea{Multi-source BFS starting from every gate.}
\complexity{O(m * n) time and space}
\begin{code}
class Solution:
    def wallsAndGates(self, rooms: List[List[int]]) -> None:
        INF = 2147483647
        R, C = len(rooms), len(rooms[0])
        q = deque((r, c) for r in range(R) for c in range(C) if rooms[r][c] == 0)
        while q:
            r, c = q.popleft()
            for nr, nc in ((r + 1, c), (r - 1, c), (r, c + 1), (r, c - 1)):
                if 0 <= nr < R and 0 <= nc < C and rooms[nr][nc] == INF:
                    rooms[nr][nc] = rooms[r][c] + 1
                    q.append((nr, nc))
\end{code}
\end{problem}

\begin{problem}{207}{Course Schedule}{Medium}
\summary{Given prerequisite pairs, can all courses be finished (i.e. no cycle)?}
\keyidea{Kahn's topological sort: repeatedly take in-degree-0 nodes; success if all are taken.}
\complexity{O(V + E) time and space}
\begin{code}
class Solution:
    def canFinish(self, numCourses: int, prerequisites: List[List[int]]) -> bool:
        adj = [[] for _ in range(numCourses)]
        indeg = [0] * numCourses
        for a, b in prerequisites:          # b -> a
            adj[b].append(a)
            indeg[a] += 1
        q = deque(i for i in range(numCourses) if indeg[i] == 0)
        taken = 0
        while q:
            c = q.popleft()
            taken += 1
            for nxt in adj[c]:
                indeg[nxt] -= 1
                if indeg[nxt] == 0:
                    q.append(nxt)
        return taken == numCourses
\end{code}
\end{problem}

\begin{problem}{210}{Course Schedule II}{Medium}
\summary{Return a valid order to take all courses, or [] if impossible.}
\keyidea{Same Kahn's algorithm; the pop order is the answer.}
\complexity{O(V + E) time and space}
\begin{code}
class Solution:
    def findOrder(self, numCourses: int, prerequisites: List[List[int]]) -> List[int]:
        adj = [[] for _ in range(numCourses)]
        indeg = [0] * numCourses
        for a, b in prerequisites:
            adj[b].append(a)
            indeg[a] += 1
        q = deque(i for i in range(numCourses) if indeg[i] == 0)
        order = []
        while q:
            c = q.popleft()
            order.append(c)
            for nxt in adj[c]:
                indeg[nxt] -= 1
                if indeg[nxt] == 0:
                    q.append(nxt)
        return order if len(order) == numCourses else []
\end{code}
\end{problem}

\begin{problem}{684}{Redundant Connection}{Medium}
\summary{A tree plus one extra edge; return the edge (latest in input) that creates the cycle.}
\keyidea{Union-Find: the first edge whose endpoints already share a root is redundant.}
\complexity{O(n * alpha(n)) time, O(n) space}
\begin{code}
class Solution:
    def findRedundantConnection(self, edges: List[List[int]]) -> List[int]:
        parent = list(range(len(edges) + 1))
        def find(x):
            while parent[x] != x:
                parent[x] = parent[parent[x]]   # path compression
                x = parent[x]
            return x
        for a, b in edges:
            ra, rb = find(a), find(b)
            if ra == rb:
                return [a, b]
            parent[ra] = rb
\end{code}
\end{problem}

\begin{problem}{323}{Number of Connected Components in an Undirected Graph}{Medium}
\summary{n nodes and an edge list; count connected components.}
\keyidea{Union-Find: start with n components, subtract 1 for every successful union.}
\complexity{O(E * alpha(n)) time, O(n) space}
\begin{code}
class Solution:
    def countComponents(self, n: int, edges: List[List[int]]) -> int:
        parent = list(range(n))
        def find(x):
            while parent[x] != x:
                parent[x] = parent[parent[x]]
                x = parent[x]
            return x
        comps = n
        for a, b in edges:
            ra, rb = find(a), find(b)
            if ra != rb:
                parent[ra] = rb
                comps -= 1
        return comps
\end{code}
\end{problem}

\begin{problem}{261}{Graph Valid Tree}{Medium}
\summary{Do n nodes and the given undirected edges form a valid tree?}
\keyidea{A tree has exactly n-1 edges and no cycle (equivalently: n-1 edges and connected).}
\complexity{O(n * alpha(n)) time, O(n) space}
\begin{code}
class Solution:
    def validTree(self, n: int, edges: List[List[int]]) -> bool:
        if len(edges) != n - 1:
            return False
        parent = list(range(n))
        def find(x):
            while parent[x] != x:
                parent[x] = parent[parent[x]]
                x = parent[x]
            return x
        for a, b in edges:
            ra, rb = find(a), find(b)
            if ra == rb:
                return False
            parent[ra] = rb
        return True
\end{code}
\end{problem}

\begin{problem}{127}{Word Ladder}{Hard}
\summary{Fewest words in a transformation from beginWord to endWord, changing one letter at a time via wordList.}
\keyidea{BFS; neighbours found via wildcard patterns like 'h*t' -\textgreater{} words.}
\complexity{O(N * L\textasciicircum{}2) time, O(N * L\textasciicircum{}2) space}
\begin{code}
class Solution:
    def ladderLength(self, beginWord: str, endWord: str, wordList: List[str]) -> int:
        if endWord not in wordList:
            return 0
        patterns = defaultdict(list)
        for w in wordList + [beginWord]:
            for i in range(len(w)):
                patterns[w[:i] + "*" + w[i + 1:]].append(w)
        seen, q, steps = {beginWord}, deque([beginWord]), 1
        while q:
            for _ in range(len(q)):
                w = q.popleft()
                if w == endWord:
                    return steps
                for i in range(len(w)):
                    for nb in patterns[w[:i] + "*" + w[i + 1:]]:
                        if nb not in seen:
                            seen.add(nb)
                            q.append(nb)
            steps += 1
        return 0
\end{code}
\end{problem}
